\chapter*{\centerline{DAFTAR ISTILAH DAN SINGKATAN}}
\addcontentsline{toc}{chapter}{DAFTAR ISTILAH DAN SINGKATAN}

\noindent\textbf{Singkatan}

{\footnotesize
\setlength{\tabcolsep}{4pt}
\renewcommand{\arraystretch}{1.05}
\begin{longtable}{|L{0.17\textwidth}|L{0.72\textwidth}|}
\hline
\textbf{Singkatan} & \textbf{Kepanjangan atau arti} \\ \hline
\endfirsthead
\hline
\textbf{Singkatan} & \textbf{Kepanjangan atau arti} \\ \hline
\endhead
Acc@$k$ & Akurasi \textit{top-k}, yaitu proporsi \textit{record} yang kode pada \textit{gold standard} termasuk dalam $k$ keluaran teratas \\ \hline
ADR & \textit{Action Design Research} \\ \hline
API & \textit{Application Programming Interface} (antarmuka pemrograman aplikasi) \\ \hline
B0, B1, B1$'$ & Kondisi pembanding eksperimen standardisasi: \textit{baseline} leksikal, LLM tanpa \textit{retrieval}, dan LLM dengan daftar golongan pokok di konteks \\ \hline
BPMN & \textit{Business Process Model and Notation} \\ \hline
DKUKMPP & Dinas Koperasi, Usaha Kecil dan Menengah, Perdagangan dan Perindustrian \\ \hline
DSR & \textit{Design Science Research} \\ \hline
DSRM & \textit{Design Science Research Methodology} \\ \hline
FEDS & \textit{Framework for Evaluation in Design Science} \\ \hline
GLMM & \textit{Generalized Linear Mixed Model} \\ \hline
GPU & \textit{Graphics Processing Unit} \\ \hline
HR & \textit{Hallucination rate}, tingkat halusinasi kode \\ \hline
JCIS & \textit{Journal of Computer Information Systems} \\ \hline
KBLI & Klasifikasi Baku Lapangan Usaha Indonesia \\ \hline
KF, KNF & Kebutuhan fungsional, kebutuhan nonfungsional \\ \hline
LKjIP & Laporan Kinerja Instansi Pemerintah \\ \hline
LLM & \textit{Large Language Model} \\ \hline
NIB & Nomor Induk Berusaha \\ \hline
NIK & Nomor Induk Kependudukan \\ \hline
P1, P2a, P2b & Kondisi uji eksperimen standardisasi: \textit{retrieval-grounded} LLM, ditambah glosarium istilah lokal, dan ditambah contoh \textit{few-shot} \\ \hline
PerBPS & Peraturan Badan Pusat Statistik \\ \hline
PP & Peraturan Pemerintah \\ \hline
RAG & \textit{Retrieval-Augmented Generation} \\ \hline
Renstra & Rencana Strategis \\ \hline
RQ & Pertanyaan penelitian (\textit{research question}) \\ \hline
RTLX & \textit{Raw NASA Task Load Index}, skor beban kerja tanpa pembobotan berpasangan \\ \hline
SAW & \textit{Simple Additive Weighting} \\ \hline
SPK & Sistem pendukung keputusan \\ \hline
SUS & \textit{System Usability Scale} \\ \hline
UMKM & Usaha mikro, kecil, dan menengah \\ \hline
UU PDP & Undang-Undang Pelindungan Data Pribadi \\ \hline
\end{longtable}
}

\vspace{1em}
\noindent\textbf{Istilah}

{\footnotesize
\setlength{\tabcolsep}{4pt}
\renewcommand{\arraystretch}{1.05}
\begin{longtable}{|L{0.27\textwidth}|L{0.62\textwidth}|}
\hline
\textbf{Istilah} & \textbf{Penjelasan} \\ \hline
\endfirsthead
\hline
\textbf{Istilah} & \textbf{Penjelasan} \\ \hline
\endhead
\textit{Audit trail} & Catatan jejak seluruh perubahan data dan keputusan verifikasi beserta pelakunya \\ \hline
\textit{Baseline} & Kondisi pembanding paling sederhana yang dipakai sebagai acuan kinerja \\ \hline
\textit{Benchmark} & Himpunan data atau tugas baku untuk membandingkan kinerja model \\ \hline
\textit{Cluster bootstrap} & Pengambilan sampel ulang pada tingkat kelompok (di sini desa) untuk memperhitungkan ketergantungan antar-\textit{record} \\ \hline
\textit{Confusion matrix} & Tabel yang membandingkan kelas pada \textit{gold standard} dengan kelas prediksi \\ \hline
\textit{Coverage--accuracy curve} & Kurva yang memperlihatkan akurasi selektif pada berbagai tingkat cakupan keluaran yang diterima otomatis \\ \hline
\textit{Data profiling} & Pemeriksaan statistik awal terhadap isi data, misalnya jumlah \textit{record}, nilai kosong, dan sebaran nilai \\ \hline
\textit{Design effect} & Faktor pengali ukuran sampel akibat ketergantungan antar-observasi dalam satu kelompok \\ \hline
\textit{Embedding} & Representasi vektor teks yang dipakai untuk mengukur kemiripan makna \\ \hline
\textit{Few-shot} & Pemberian beberapa contoh pada \textit{prompt} agar model memahami tugas \\ \hline
\textit{Gold standard} & Label acuan hasil pengodean manusia yang dipakai untuk menilai keluaran model \\ \hline
\textit{Inter-annotator agreement} & Tingkat kesepakatan antaranotator yang mengodekan data yang sama \\ \hline
\textit{Open-weight} & Model yang bobotnya tersedia sehingga dapat dijalankan pada infrastruktur sendiri \\ \hline
\textit{Output reproducibility} & Keterulangan keluaran, yaitu proporsi \textit{record} yang memperoleh keluaran identik pada beberapa kali pengulangan \\ \hline
\textit{Pipeline} & Rangkaian langkah pemrosesan data dari masukan hingga keluaran \\ \hline
\textit{Prompt} & Instruksi dan konteks yang diberikan kepada LLM \\ \hline
\textit{Proprietary} & Model yang hanya dapat diakses melalui layanan penyedia, biasanya melalui API \\ \hline
\textit{Prototype} & Versi awal sistem yang dibangun sebagai intervensi sekaligus instrumen penelitian \\ \hline
\textit{Record} & Satu baris data, yaitu satu entri usaha \\ \hline
\textit{Re-ranking} & Pengurutan ulang kandidat hasil \textit{retrieval} \\ \hline
\textit{Retrieval-grounded} & Keluaran model dibatasi dan dipandu oleh kandidat yang diambil dari sumber acuan \\ \hline
\textit{Retriever} & Komponen yang mengambil kandidat relevan dari basis pengetahuan \\ \hline
\textit{Threshold} & Nilai ambang yang memisahkan keluaran yang diterima otomatis dari yang diverifikasi manusia \\ \hline
\textit{Usability} & Tingkat kemudahan, efisiensi, dan kepuasan pengguna dalam memakai sistem \\ \hline
\textit{Within-subject} & Rancangan ketika setiap partisipan mengalami semua kondisi \\ \hline
\textit{Wilcoxon signed-rank test} & Uji nonparametrik untuk selisih berpasangan \\ \hline
\textit{Zero-shot} & Penggunaan model tanpa contoh pada \textit{prompt} \\ \hline
\end{longtable}
}

\clearpage
